\documentclass[../../main.tex]{subfiles}

\begin{document}
\chapter{Manual de usuario}

% Introducción 
% Instalación
% Interfaz
% Proyectos
% Edición
% Búsqueda
% Compilación
% Terminal
% Configuración
% Atajos de Teclado
% FAQ

%% Seccion de Introducción    
\section{Introducción}
\textbf{SamaruC Retro Editor} es un entorno de desarrollo integrado (IDE) especializado para programación en lenguaje C, diseñado con una interfaz moderna y funcional basada en JavaFX.

\par{\textbf{Nota:} El nombre \textit{SamaruC} es un juego de palabras. Por una parte \textit{Samaruc} es un pez endémico de la Comunidad Valenciana en grave peligro de extinción, referenciando la región de origen del proyecto. Por otra parte, remarca la terminación con letra mayuscula de la \texttt{C}, haciendo referencia al lenguaje de programación para el cual está optimizado este editor.}

\subsection{Características principales}
\begin{itemize}
    \item \textbf{Editor de código} con resaltado de sintaxis para C/C++
    \item \textbf{Gestión de proyectos} con explorador de archivos integrado
    \item \textbf{Compilación y ejecución} directa desde el IDE
    \item \textbf{Terminal integrada} con consola de salida y depuración
    \item \textbf{Búsqueda avanzada} en archivos y proyectos completos
    \item \textbf{Pestañas múltiples} para trabajar con varios archivos simultáneamente
    \item \textbf{Interfaz personalizable} con temas claro y oscuro
    \item \textbf{Soporte multiidioma} (Español/Inglés)
    \item \textbf{Deshacer/Rehacer} con historial ilimitado
    \item \textbf{Iconos FontAwesome 5} para una interfaz moderna
\end{itemize}

\clearpage

%% Seccion de instalación
\section{Instalación}
\subsection{Requisitos del Sistema}
\begin{itemize}
    \item \textbf{Java:} JDK 17 o superior
    \item \textbf{Maven:} 3.6 o superior (para compilar desde código fuente)
    \item \textbf{Compilador C:} GCC, Z88DK o GBDK (según tu proyecto)
    \item \textbf{Sistema Operativo:} Windows, Linux o macOS
\end{itemize}

\subsection{Instalación desde Código Fuente}
Para compilar y ejecutar el editor desde el código fuente:
\begin{enumerate}
    \item Clona o descarga el repositorio del proyecto
    \item Abre una terminal en el directorio del proyecto
    \item Ejecuta los siguientes comandos:
\end{enumerate}

\begin{lstlisting}[language = Java]
mvn clean compile
mvn install
mvn javafx:run
\end{lstlisting}

\subsection{Instalación desde Ejecutable}
Si tienes el paquete ejecutable (.exe en Windows):
\begin{itemize}
    \item Extrae el archivo \texttt{RetroEditor-win.zip}
    \item Ejecuta \texttt{RetroEditor.exe}
    \item No requiere instalación de Java (incluye runtime embebido)
\end{itemize}

\clearpage

%sección de interfaz
\section{Interfaz}
\subsection{Descripción de la Interfaz}

\subsubsection{Barra de Menús}
La barra de menús se encuentra en la parte superior de la ventana como se puede apreciar en la captura y contiene los siguientes menús principales:

\begin{figure}[!htb]
    \centering
    \includegraphics[width=10cm]{aplicacion/Menus.png}
    \caption{Situación de los menus desplegables de la herramienta.}
    \label{fig:Menus}
\end{figure}

\par{ Funciones principales de cada menú: }
\begin{itemize}
    \item \textbf{Archivo:} Crear, abrir, guardar y cerrar archivos. 
    \item \textbf{Edición:} Deshacer, rehacer, cortar, copiar, pegar. 
    \item \textbf{Compilación:} Compilar y ejecutar el proyecto actual.
    \item \textbf{Configuración:} Ajustes de idioma, compilador y preferencias.
    \item \textbf{Ayuda:} Acceso a este manual y acerca de.
\end{itemize}

\clearpage

\subsubsection{Barra de Herramientas}
Como se puede ver en la siguiente captura de pantalla otra manera de acceder rápido a las funciones elementales de la herramienta es mediante la barra de herramientas con iconos representativos de cada función:

\begin{figure}[!htb]
    \centering
    \includegraphics[width=10cm]{aplicacion/Captura_barra_herramientas.png }
    \caption{Barra de herramientas con los iconos representativos de cada función.}
    \label{fig:BarraHerramientas}
\end{figure}

\par{La barra de herramientas proporciona acceso rápido a las funciones más utilizadas:}

\begin{itemize}

    \item \faFolderPlus\ \textbf{Nuevo Proyecto:} Abre el dialogo de creación de un nuevo proyecto
    \item \faFolderOpen\ \textbf{Abrir Proyecto:} Abre un proyecto o archivo existente
    \item \faFile\ \textbf{Nuevo:} Crea un archivo nuevo
    \item \faSave\ \textbf{Guardar:} Guardar el archivo actual
    \item \faTimes\ \textbf{Cerrar:} Cierra el archivo actual
    
    \item \faUndo\ \textbf{Deshacer:} Desacer la última acción realizada
    \item \faRedo\ \textbf{Rehacer:} Rehacer la última acción deshecha
    \item \faSearch\ \textbf{Buscar en archivo:} Herramientas de búsqueda en el archivo actual
    \item \faSearchPlus\ \textbf{Buscar en proyecto:} Herramientas de búsqueda en todo el proyecto

    \item \faCut\ \textbf{Cortar:} Función para cortar un texto seleccionado
    \item \faCopy\ \textbf{Copiar:} Función para copiar un texto seleccionado
    \item \faPaste\ \textbf{Pegar:} Función de pegado de texto copiado o cortado
    
    \item \faHammer\ \textbf{Compilar:} Construcción y compilación del código fuente
    \item \faPlay\ \textbf{Ejecutar:} Construcción, compilación y ejecución del código fuente

    \item \faCog\ \textbf{Configuración:} Acceso a preferencias y configuración
    \item \faAdjust\ \textbf{Tema:} Alternar entre tema claro y oscuro
    
    \item \faQuestionCircle\ \textbf{Ayuda - Manual:} Abre este manual de usuario
\end{itemize}

\clearpage

\subsubsection{Panel del Explorador de Proyectos}


\par{Ubicado en el lado izquierdo de la interfaz, como se aprecia en la imagen, este panel muestra la estructura de directorios y archivos del proyecto activo.}



\begin{figure}[!htb]
    \centering
    \includegraphics[width=7cm]{aplicacion/Area_explorador.jpg}
    \caption{Panel del explorador de proyectos.}
    \label{fig:ExploradorProyectos}
\end{figure}
\par{Desde este apartado es posible navegar por la jerarquía de carpetas, así como desplegar o contraer los directorios para visualizar u ocultar su contenido. Asimismo, los archivos pueden abrirse mediante doble clic, lo que permite visualizar y editar su contenido en el área de trabajo principal.}

\par{Al hacer clic con el botón derecho del ratón sobre esta zona, se despliega un menú contextual que ofrece opciones como la creación de un nuevo archivo o una nueva carpeta.}

\par{Si pulsamos la tecla \texttt{Suprimir} con un archivo o carpeta seleccionada, aparecerá un diálogo de confirmación que si confirmamos eliminará el elemento seleccionado.}

\par{Además, tenemos un botón en la esquina superior izquierda del panel del explorador que permite refrescar la vista para mostrar los cambios realizados en el sistema de archivos, como la adición o eliminación de archivos fuera del editor.}

\clearpage

\subsubsection{Área de Edición con Pestañas}
El área central donde se editan los archivos, como se puede ver en la captura. Características:

\begin{figure}[!htb]
    \centering
    \includegraphics[width=10cm]{aplicacion/CapturaZonaEdicion.png}
    \caption{Área de edición con pestañas.}
    \label{fig:AreaEdicion}
\end{figure}

\begin{itemize}
    \item Soporte para múltiples archivos abiertos simultáneamente
    \item Resaltado de sintaxis automático para C/C++
    \item Numeración de líneas
    \item Indicador de archivo modificado (*) en la pestaña
\end{itemize}

 \subsubsection{Panel de Consola Y Monitor + Barra de Estado}

 \begin{figure}[!htb]
    \centering
    \includegraphics[width=16cm]{aplicacion/ConsolaCaptura.jpg}
    \caption{Panel de consola y monitor.}
    \label{fig:PanelConsola}
\end{figure}

\begin{figure}[!htb]
    \centering
    \includegraphics[width=16cm]{aplicacion/BarraEstadoCaptura.jpg}
    \caption{Barra de estado.}
    \label{fig:BarraEstado}
\end{figure}

 Este panel se divide en dos pestañas principales: \textbf{Consola} y \textbf{Monitor}, cada una con funciones específicas para la salida de compilación, ejecución y monitorización de la aplicación en tiempo real. Además, en la parte inferior de la ventana se encuentra la barra de estado que muestra el estado actual del editor.

\clearpage

\section{Gestión de Proyectos}

En esta sección se describen las herramientas para crear, abrir y organizar proyectos de programación en SamaruC Retro Editor.

SamaruC Retro Editor utiliza una estructura de proyecto basada en directorios, donde cada proyecto corresponde a una carpeta que contiene los archivos fuente (.c), encabezados (.h) y otros recursos relacionados. El panel del explorador de proyectos permite gestionar esta estructura de manera visual e intuitiva.

\subsection{Crear un Nuevo Proyecto}

\begin{figure}[!htb]
    \centering
    \includegraphics[width=12cm]{manual/Captura_manual_nuevo_proyecto.jpg}
    \caption{Crear un nuevo proyecto.}
    \label{fig:CrearNuevoProyecto}
\end{figure}

\begin{enumerate}
    \item Haz clic en el botón \textbf{Nuevo Proyecto} o usa el menú \textit{Archivo → Nuevo Proyecto}
    \item Selecciona o crea una carpeta para el proyecto
    \item El explorador mostrará la estructura del directorio
    \item Comienza a crear archivos .c y .h
\end{enumerate}

\clearpage 

\subsection{Abrir un Proyecto Existente}

\begin{figure}[!htb]
    \centering
    \includegraphics[width=12cm]{manual/Captura_manual_abrir_proyecto.jpg}
    \caption{Abrir un proyecto existente.}
    \label{fig:AbrirProyectoExistente}
\end{figure}

\begin{enumerate}
    \item Haz clic en \textbf{Abrir Proyecto} o usa el menú \textit{Archivo → Abrir Proyecto}
    \item Navega hasta la carpeta del proyecto
    \item Selecciona el directorio raíz del proyecto
    \item Todos los archivos del proyecto aparecerán en el explorador
 \end{enumerate}

 \textbf{Consejo:} Organiza tus proyectos con subdirectorios (src/, include/, build/) para mejor estructura y mantenimiento.

 \subsection{Trabajar con Archivos}

 SamaruC Retro Editor permite crear, abrir, guardar y cerrar archivos de manera sencilla desde el menú de archivos o la barra de herramientas. Las operaciones básicas incluyen:
 
 \begin{itemize}
     \item \textbf{Crear archivo nuevo:} Botón "Nuevo" o \texttt{Ctrl+N}
     \item \textbf{Abrir archivo:} Doble clic en el explorador o botón "Abrir"
     \item \textbf{Guardar archivo:} \texttt{Ctrl+S} o botón "Guardar"
     \item \textbf{Guardar como:} Menú \textit{Archivo → Guardar Como}
     \item \textbf{Cerrar archivo:} Botón "Cerrar" o click en X de la pestaña
 \end{itemize}

\clearpage 

\section{Edición de Código}

En esta sección se describen las herramientas y características disponibles para la edición de código fuente en SamaruC Retro Editor, incluyendo el resaltado de sintaxis, operaciones de edición y características avanzadas del editor.

\subsection{Resaltado de Sintaxis}

El editor de texto integrado en SamaruC Retro Editor utiliza la biblioteca RichTextFX para proporcionar resaltado de sintaxis automático para el lenguaje C. Esto facilita la lectura y escritura de código al diferenciar visualmente los distintos elementos del lenguaje como se puede ver en la siguiente captura de pantalla.

\begin{figure}[!h]
    \centering
    \includegraphics[width=8cm]{manual/Captura_manual_ejemplo_resaltado_sintaxis.jpg}
    \caption{Resaltado de sintaxis en el editor de código.}
    \label{fig:ResaltadoSintaxis}
\end{figure}



El editor proporciona resaltado automático para las siguientes categorías de elementos de código:

\begin{itemize}
     \item \textcolor{blue}{\textbf{Palabras clave:}} if, else, for, while, return, struct, etc.
     \item \textcolor{green}{\textbf{Comentarios:}} // comentarios de línea y /* bloques */
     \item \textcolor{red}{\textbf{Cadenas:}} "texto entre comillas"
     \item \textcolor{teal}{\textbf{Números:}} 123, 0xFF, 3.14
     \item \textcolor{brown}{\textbf{Funciones:}} printf, scanf, main
     \item \textcolor{blue}{\textbf{Directivas:}} \#include, \#define, \#ifdef
\end{itemize}

\clearpage 

\subsection{Operaciones de Edición}

En la siguiente tabla se resumen las operaciones de edición más comunes y sus atajos de teclado correspondientes para mejorar la eficiencia al escribir código:

\begin{tabular}{|l|l|l|}
    \hline
    \textbf{Operación} & \textbf{Atajo} & \textbf{Descripción} \\
    \hline
    Deshacer & \texttt{Ctrl+Z} & Revierte la última acción \\
    \hline
    Rehacer & \texttt{Ctrl+Y} & Restaura la acción deshecha \\
    \hline
    Cortar & \texttt{Ctrl+X} & Corta el texto seleccionado \\
    \hline
    Copiar & \texttt{Ctrl+C} & Copia el texto seleccionado \\
    \hline
    Pegar & \texttt{Ctrl+V} & Pega desde el portapapeles \\
    \hline
    Seleccionar todo & \texttt{Ctrl+A} & Selecciona todo el texto \\
    \hline
\end{tabular}

\subsection{Características Avanzadas}

SamaruC Retro Editor incluye varias características avanzadas para mejorar la experiencia de edición de código y la productividad del programador:
\begin{itemize}
     \item \textbf{Numeración de líneas:} Siempre visible para referencia
     \item \textbf{Múltiples pestañas:} Trabaja con varios archivos simultáneamente
     \item \textbf{Detección de cambios:} Asterisco (*) indica archivos modificados no guardados
     \item \textbf{Historial ilimitado:} Deshacer/Rehacer sin límites
\end{itemize}

\clearpage

\section{Búsqueda de Texto}

SamaruC Retro Editor ofrece potentes herramientas de búsqueda para localizar texto dentro de un archivo o en todo el proyecto. Estas funciones son esenciales para navegar rápidamente por el código, encontrar definiciones de funciones, referencias a variables o cualquier término específico.

\subsection{Buscar en Archivo Actual}

Cuando deseas buscar un término dentro del archivo que estás editando, puedes utilizar la función de búsqueda integrada de la siguiente manera:


\begin{itemize}
    \item Presiona \texttt{Ctrl+F} o haz clic en el botón de búsqueda       
    \item Se abre un panel de búsqueda en la parte superior del editor
    \item Escribe el término a buscar
    \item Navega entre resultados con los botones Anterior/Siguiente
\end{itemize}

\clearpage

\subsection{Buscar en Todo el Proyecto}

Si necesitas encontrar un término en todos los archivos de tu proyecto, la función de búsqueda en proyecto te permite hacerlo de manera eficiente:

\begin{figure}[!h]
    \centering
    \includegraphics[width=8cm]{manual/Captura_manual_buscar_en_proyecto.jpg}
    \caption{Búsqueda en proyecto actual.}
    \label{fig:BuscarEnProyecto}
\end{figure}

\begin{itemize}
    \item Presiona \texttt{Ctrl+H} o usa el botón  en la barra de herramientas
    \item Busca en todos los archivos del proyecto
    \item Los resultados muestran archivo, línea y contexto
    \item Haz doble clic en un resultado para abrir el archivo en esa ubicación
\end{itemize}

\begin{figure}[!h]
    \centering
    \includegraphics[width=6cm]{manual/Captura_manual_buscar_en_proyecto_resultados.jpg}
    \caption{Resultados de búsqueda en proyecto.}
    \label{fig:ResultadosBusquedaProyecto}
\end{figure}

\textbf{Consejo:} Usa la búsqueda en proyecto para encontrar definiciones de funciones, referencias a variables o identificar dónde se usa una función específica.

\clearpage

% %% TODO a partir de aquí, revisar y completar el contenido de cada sección 

%         \section{Compilación y Ejecución}

% \subsection{Compilar el Proyecto}
% \begin{enumerate}
%     \item Asegúrate de tener un proyecto abierto con archivos .c
%     \item Haz clic en el botón \textbf{Compilar} o usa el menú \textit{Compilación → Compilar}
%     \item La salida de compilación aparece en la consola inferior
%     \item Los errores y warnings se muestran con detalles de línea y archivo
% \end{enumerate}

% \subsection{Ejecutar el Programa}
% \begin{enumerate}
%     \item Después de compilar exitosamente, haz clic en \textbf{Ejecutar}
%     \item El programa se ejecuta y la salida aparece en la consola
%     \item La entrada interactiva también está disponible en la consola
% \end{enumerate}

% \subsection{Configurar el Compilador}
% El editor soporta diferentes compiladores C:
% \begin{itemize}
%     \item \textbf{GCC:} Compilador GNU para C/C++
%     \item \textbf{MinGW:} GCC para Windows
%     \item \textbf{GBDK:} Game Boy Development Kit
% \end{itemize}

% Configura la ruta del compilador en \textit{Configuración → Compilador}

% \textbf{Nota:} Asegúrate de que el compilador esté en el PATH del sistema o configura la ruta completa en las preferencias.

%         \section{Terminal y Consola}

% \subsection{Consola de Salida}
% Muestra:
% \begin{itemize}
%     \item Salida del compilador (gcc, make, etc.)
%     \item Mensajes de error y warnings
%     \item Salida del programa (stdout)
%     \item Entrada interactiva (stdin)
% \end{itemize}

% \subsection{Consola de Depuración}
% Registra:
% \begin{itemize}
%     \item Mensajes de log del sistema
%     \item Información de depuración del IDE
%     \item Errores internos y diagnósticos
% \end{itemize}

% \subsection{Limpiar la Consola}
% Usa el botón de limpiar en el panel de consola para borrar la salida anterior y tener una vista limpia.

%         \section{Configuración y Personalización}

% \subsection{Configuración de Idioma}
% \begin{enumerate}
%     \item Abre \textit{Configuración → Preferencias} o haz clic en el botón de configuración
%     \item Ve a la pestaña \textbf{Idioma}
%     \item Selecciona entre Español o Inglés
%     \item Haz clic en \textbf{Guardar}
%     \item La interfaz se actualiza automáticamente
% \end{enumerate}

% \subsection{Configuración del Compilador}
% \begin{enumerate}
%     \item Abre las preferencias () % aqui el icono de engranaje %% TODO
%     \item Ve a la pestaña \textbf{Compilador}
%     \item Especifica la ruta completa al ejecutable del compilador
%     \item Ejemplo: \texttt{C:\textbackslash MinGW\textbackslash bin\textbackslash gcc.exe} o \texttt{/usr/bin/gcc}
%     \item Guarda los cambios
% \end{enumerate}

% \subsection{Temas Visuales}
% El editor incluye dos temas:
% \begin{itemize}
%     \item \textbf{Tema Claro:} Fondo blanco, ideal para ambientes iluminados
%     \item \textbf{Tema Oscuro:} Fondo oscuro, reduce fatiga visual en ambientes con poca luz
% \end{itemize}

% Alterna entre temas usando el botón en la barra de herramientas.

% \subsection{Archivo de Configuración}
% Las preferencias se guardan en: \texttt{\~{}/.retroeditor.properties}

% Puedes editar manualmente este archivo si es necesario.





 \clearpage
 %% Seccion de miscelánea
\section{Miscelánea}

\subsection{Atajos de Teclado}

\par{Para mejorar la eficiencia y la experiencia de usuario, el editor incluye una serie de atajos de teclado para las funciones más comunes. Estos atajos permiten realizar acciones rápidamente sin necesidad de navegar por los menús, lo que agiliza el flujo de trabajo. A continuación se presenta una tabla con los atajos de teclado disponibles en \textbf{SamaruC Retro Editor}:}
\begin{table}[ht]
    \centering
    \small
    \begin{tabular}{p{3cm}p{6cm}}
        \toprule
            \textbf{Atajo} & \textbf{Acción} \\
        \midrule
            \texttt{Ctrl+N} & Nuevo archivo \\
            \texttt{Ctrl+O} & Abrir archivo \\
            \texttt{Ctrl+S} & Guardar archivo \\
            \texttt{Ctrl+W} & Cerrar archivo actual \\
            \texttt{Ctrl+Z} & Deshacer \\
            \texttt{Ctrl+Y} & Rehacer \\
            \texttt{Ctrl+X} & Cortar \\
            \texttt{Ctrl+C} & Copiar \\
            \texttt{Ctrl+V} & Pegar \\
            \texttt{Ctrl+A} & Seleccionar todo \\
            \texttt{Ctrl+F} & Buscar en archivo actual \\
            \texttt{Ctrl+H} & Buscar en proyecto completo \\
            \texttt{F5} & Compilar proyecto \\
            \texttt{F6} & Ejecutar programa \\
        \bottomrule
    \end{tabular}
    \caption{Atajos de teclado disponibles en SamaruC Retro Editor.}
    \label{tab:atajos_teclado}
\end{table}

\subsection{Tecnologías utilizadas en el desarrollo de SamaruC Retro Editor:}
\begin{itemize}
    \item \textbf{JavaFX 17:} Framework de interfaz de usuario
    \item \textbf{RichTextFX:} Editor de texto con resaltado de sintaxis
    \item \textbf{Ikonli FontAwesome 5:} Iconos de interfaz
    \item \textbf{Maven:} Gestión de dependencias y construcción
\end{itemize}

\clearpage

\subsection{Requisitos del Sistema:}
\begin{itemize}
    \item Java 17+
    \item 4 GB RAM mínimo (8 GB recomendado)
    \item 200 MB espacio en disco
    \item Compilador C (Z88DK o GBDK)
\end{itemize}

\end{document}